\documentclass[11pt,a4paper]{article}

\usepackage[margin=2.6cm]{geometry}
\usepackage{amsmath,amssymb,amsthm,mathrsfs,mathtools}
\usepackage{xcolor}
\usepackage{dsfont}
\usepackage[T1]{fontenc}
\usepackage{microtype}
\usepackage{booktabs}
\usepackage[colorlinks=true,linkcolor=blue!55!black,citecolor=blue!55!black,urlcolor=blue!55!black]{hyperref}

\newtheorem{thm}{Theorem}[section]
\newtheorem{prop}[thm]{Proposition}
\newtheorem{lemma}[thm]{Lemma}
\newtheorem{cor}[thm]{Corollary}
\theoremstyle{definition}
\newtheorem{defn}[thm]{Definition}
\newtheorem{rem}[thm]{Remark}

\newcommand{\fA}{\mathfrak{A}}
\newcommand{\fB}{\mathfrak{B}}
\newcommand{\fh}{\mathfrak{h}}
\newcommand{\cH}{\mathcal{H}}
\newcommand{\cD}{\mathcal{D}}
\newcommand{\vep}{\varepsilon}
\newcommand{\1}{\mathds{1}}
\newcommand{\CAR}{\textup{CAR}}
\newcommand{\sign}{\operatorname{sign}}
\newcommand{\ran}{\operatorname{ran}}
\newcommand{\rank}{\operatorname{rank}}
\newcommand{\tr}{\operatorname{tr}}
\newcommand{\ltwo}{\ell^{2}}
\newcommand{\R}{\mathbb{R}}
\newcommand{\C}{\mathbb{C}}
\newcommand{\Z}{\mathbb{Z}}
\newcommand{\N}{\mathbb{N}}
\newcommand{\Gam}{\Gamma}

% numbers of the companion manuscript, generated by sync_docs.py
% AUTO-GENERATED by sync_docs.py from build/free_fermion_cft_v5.aux -- do not edit.
% Numbers of the companion manuscript free_fermion_cft_v5.tex.
\newcommand{\mrRGiso}{87}
\newcommand{\mrRGflow}{89}
\newcommand{\mrSL}{90}
\newcommand{\mrChiral}{61}
\newcommand{\mrLatVac}{60}
\newcommand{\mrHardy}{129}
\newcommand{\mrMomRG}{121}
\newcommand{\mrAlias}{249}
\newcommand{\mrHS}{25}
\newcommand{\mrDefMomRG}{3.16}
\newcommand{\mrDefWavRG}{3.2}
\newcommand{\mrJackson}{4.28}
\newcommand{\mrMoeb}{4.22}
\newcommand{\mrThmC}{4.16}
\newcommand{\mrThmSm}{4.25}


\title{\bf Making the Zini--Wang comparison precise\\[0.3cm]
\large A note on \S1.2 of \emph{Conformal field theory from lattice fermions}}
\author{}
\date{31 July 2026}

\begin{document}
\maketitle

\begin{abstract}
Section 1.2 of the published article, and of the companion manuscript
\texttt{free\_fermion\_cft\_v5.tex} as it stood when this note was written, compares the
operator-algebraic renormalization (OAR) construction of the scaling limit with the
\emph{low-energy scaling limit} of Zini and Wang, and asserts that the Zini--Wang structure ``is not
immediately available'' for the massless chiral lattice fermions treated there, ``where the finitely
renormalized states $\omega^{(N)}_{M}$ are not pure (due to the entanglement of the chiral halves)''.
That assertion is stated without proof and, as it stands, identifies the wrong obstruction.\footnote{%
Every quotation of \S1.2 in this note is from the published article, Commun.\ Math.\ Phys.\
\textbf{398} (2023) 219--289, whose wording the manuscript carried unchanged until 22 September 2026.
On that date \S1.2 of \texttt{free\_fermion\_cft\_v5.tex} was rewritten in the light of this note,
and now states the multiplicity criterion of \S\ref{sec:gns} instead; the passages quoted
below are therefore no longer to be found there.} We make
the comparison precise. The renormalization flow acts on one-particle symbols by \emph{compression}
along an isometry, and we give the exact identity governing when a compression of a projection is
again a projection (Lemma \ref{lem:compression}). We then observe that what the Zini--Wang
connecting unitaries actually require is not purity but constancy of the GNS multiplicity, which for
a quasi-free state is $2^{m(S)}$ with $m(S)$ the number of eigenvalues of the symbol strictly between
$0$ and $1$ (Proposition \ref{prop:gns}). Computing $m$ in the model of the manuscript gives closed
forms in every case, and yields a sharper statement: for the chiral algebra the multiplicity is
\emph{constant in $M$}, so the connecting unitaries do exist at every finite scale, and the
obstruction is a \emph{discontinuity at $M=\infty$}, where the multiplicity collapses by a factor
$2^{|\Gam_{N,-}|}$ (Theorem \ref{thm:model}, Corollary \ref{cor:obstruction}). Quantitatively, the
chiral restriction of the finite-scale vacuum has extensive entropy with density exactly
$2\ln2-1$ per site (Corollary \ref{cor:entropy}). All closed forms are confirmed numerically in
\texttt{numerics/os\_check19.py}.
\end{abstract}

\section{Introduction and conventions}
\label{sec:intro}

\paragraph{References to the main text.}
Throughout, citations of the form \S1.2, (\mrLatVac), Definition \mrDefMomRG, Lemma \mrJackson refer to the companion
manuscript
\begin{center}
\texttt{free\_fermion\_cft\_v5.tex} \quad(the revised version of \cite{OsborneStottmeister}),
\end{center}
with the numbering of that file as compiled on 31 July 2026; the numbers are pulled from its \texttt{.aux} into \texttt{mainrefs.tex} by \texttt{sync\_docs.py}, so they stay correct when the manuscript is edited. Where a statement of the manuscript is
quoted verbatim it is set in quotation marks. Numerical checks live in
\texttt{numerics/os\_check19.py} of the same directory and are labelled \texttt{[31]}--\texttt{[34]}
there.

\paragraph{What \S1.2 asserts.}
The manuscript recalls the Zini--Wang \emph{low-energy scaling limit} \cite{ZiniWang}: a nested
family of energy-bounded subspaces $\cH^{E}_{N}\subset\cH^{E'}_{N}$, \emph{connecting unitaries}
$\phi^{E}_{N}:\cH^{E}_{N}\rightarrow\cH^{E}_{N+1}$ defined for $N$ large enough that
$\dim\cH^{E}_{N}=d_{E}$ stabilises, and the extension property
$\phi^{E}_{N}=\phi^{E'}_{N}|_{\cH^{E}_{N}}$. It then explains that OAR reproduces the Zini--Wang
compatibility condition through the GNS isometries induced by the refining morphisms, and continues:
\begin{quote}
``\dots if the GNS representations $\{\cH_{M}^{(N)},\pi_{M}^{(N)},\Omega_{M}^{(N)}\}$ of the finitely
renormalized states $\omega^{(N)}_{M}$ are unitarily equivalent, e.g.~assuming a Stone--von
Neumann-type result at finite scales, that is, if there are (connecting) unitaries
$\phi^{N}_{M}:\cH^{(N)}_{M}\rightarrow\cH^{(N)}_{M+1}$, we can recover the structure of a low-energy
scaling limit \dots such that $N$ plays the r\^ole of the energy scale $E$.''
\end{quote}
and, a page later,
\begin{quote}
``It is worth pointing out that such a structure is not immediately available for the
operator-algebraic renormalization of (free) massless chiral lattice fermions in $1+1$-dimensions
considered here, where the finitely renormalized states $\omega^{(N)}_{M}$ are not pure (due to the
entanglement of the chiral halves) in contrast with their scaling limit (given by the Hardy
projection).''
\end{quote}
Two things are left open. First, the hypothesis ``assuming a Stone--von Neumann-type result at finite
scales'' is not verified, nor is it said what it should mean. Second, the non-purity of
$\omega^{(N)}_{M}$ is asserted, not proved, and the inference from non-purity to the unavailability
of the Zini--Wang structure is not justified. Sections
\ref{sec:compression}--\ref{sec:obstruction} settle all three points. Section \ref{sec:open} records
what this note does \emph{not} settle.

\paragraph{Setting.}
We work with the OAR data of \S3 of the manuscript. At each scale $N$ we have a finite-dimensional
one-particle space $\fh_{N}$ ($\dim\fh_{N}=d_{N}<\infty$), the CAR algebra
$\fA_{N}=\fA_{\CAR}(\fh_{N})$, and isometries $R^{N}_{N+1}:\fh_{N}\rightarrow\fh_{N+1}$ forming an
inductive system (eq.~(\mrRGiso) of the manuscript), with associated Bogoliubov morphisms
$\alpha^{N}_{N+1}:\fA_{N}\rightarrow\fA_{N+1}$. The renormalization group flow on states is (\mrRGflow),
\begin{align}
\label{eq:rgflow}
\omega^{(N)}_{M} & = \omega^{(N+M)}\circ\alpha^{N}_{N+M}, & N,M & \in\N_{0},
\end{align}
a state on the \emph{fixed} algebra $\fA_{N}$, and the scaling limit is
$\omega^{(N)}_{\infty}=\lim_{M}\omega^{(N)}_{M}$ (\mrSL). Recall that a quasi-free state $\omega_{S}$
on $\fA_{\CAR}(\fh)$ is determined by its symbol $S=S^{*}\in B(\fh)$ with $0\leq S\leq\1$ via
$\omega_{S}(a^{\dagger}(\xi)a(\eta))=\langle\eta,S\xi\rangle$, and that $\omega_{S}$ is pure if and
only if $S$ is a projection \cite{Araki,EvansKawahigashi}. Since all $\omega^{(N)}$ in the manuscript
are quasi-free and the $\alpha$'s are Bogoliubov, \eqref{eq:rgflow} reads at the one-particle level
\begin{align}
\label{eq:symbolflow}
S^{(N)}_{M} & = (R^{N}_{N+M})^{*}\!\ S^{(N+M)}\!\ R^{N}_{N+M},
\end{align}
that is, \emph{the renormalization of a quasi-free state is the compression of its symbol along an
isometry}. Everything below follows from this one observation.

\begin{rem}[Two families of projections]
\label{rem:notation}
The manuscript uses $\pm$ in two unrelated ways, and both occur below. We write $p_{\pm}=\tfrac{1}{2}(\1_{2}\pm\sigma_{y})$ for the \emph{chirality} projections on the spinor index $\C^{2}$ of (\mrChiral), which are constant in the momentum $k$; and $P^{\pm}$ for the \emph{Hardy} projections onto positive respectively negative momenta, which act on the momentum variable and are trivial on the spinor index. The scaling-limit symbol (\mrHardy) happens to couple the two, $P_{0}(k)=p_{\sign(k)}$, which is precisely why the chiral scaling limit is pure.
\end{rem}

\section{Compressions of quasi-free symbols}
\label{sec:compression}

\begin{lemma}[Compression identity]
\label{lem:compression}
Let $\fh',\fh$ be Hilbert spaces, $R:\fh'\rightarrow\fh$ an isometry, and $S=S^{*}\in B(\fh)$ with
$0\leq S\leq\1$. Put $S':=R^{*}SR$ and $Q:=\1-RR^{*}$, the projection onto $(\ran R)^{\perp}$. Then
$0\leq S'\leq\1$ and
\begin{align}
\label{eq:compid}
S'-S'^{2} & = R^{*}(S-S^{2})R + X^{*}X, & X & := QSR.
\end{align}
In particular, if $S$ is a projection then $S'-S'^{2}=X^{*}X\geq0$, and the following are
equivalent:
\begin{enumerate}
\item[(i)] $S'$ is a projection, i.e.~$\omega_{S'}$ is pure;
\item[(ii)] $X=QSR=0$;
\item[(iii)] $\ran R$ is invariant under $S$.
\end{enumerate}
Moreover $\|S'-S'^{2}\|=\|X\|^{2}$, and if $R$ is the inclusion of a subspace $\fh'\subset\fh$ with
orthogonal projection $p$, then $X=(\1-p)Sp$ and $\|X\|=\|[S,p]\|$. Finally, if $S$ is \emph{strictly} between $0$ and $1$, i.e.~$0<S<\1$, then so is $S'$, and hence
\begin{align}
\label{eq:strictm}
m(S') & = \dim\fh' \qquad\text{whenever }0<S<\1,
\end{align}
with $m$ as in Proposition \ref{prop:gns}.
\end{lemma}

\begin{proof}
$0\leq S'\leq\1$ is immediate from $R^{*}R=\1$. Using $RR^{*}=\1-Q$,
\begin{align*}
S'^{2} & = R^{*}SRR^{*}SR = R^{*}S(\1-Q)SR = R^{*}S^{2}R - R^{*}SQSR,
\end{align*}
and $Q=Q^{*}=Q^{2}$ gives $R^{*}SQSR=(QSR)^{*}(QSR)=X^{*}X$; subtracting from $S'=R^{*}SR$ yields
\eqref{eq:compid}. If $S^{2}=S$ the first term drops. Then $S'$ is a projection iff $X^{*}X=0$ iff
$X=0$, and $X=QSR=0$ says exactly that $S(\ran R)\subset\ran R$. The norm identity is
$\|X^{*}X\|=\|X\|^{2}$. For the last statement take $\fh'\subset\fh$ and $R$ the inclusion, so
$RR^{*}=p$ and $X=(\1-p)Sp$; since $[S,p]=(\1-p)Sp-pS(\1-p)$ is the difference of the two
off-diagonal blocks of $S$ in the decomposition $\fh=p\fh\oplus(\1-p)\fh$, and these are adjoint to
one another, $\|[S,p]\|=\|(\1-p)Sp\|$. For the last claim, if $0<S<\1$ then for $0\neq\xi\in\fh'$ we have $\langle\xi,S'\xi\rangle=\langle R\xi,S\!\ R\xi\rangle\in(0,\|R\xi\|^{2})=(0,\|\xi\|^{2})$, because $R\xi\neq0$; so no eigenvalue of $S'$ equals $0$ or $1$, which is \eqref{eq:strictm}.
\end{proof}

Lemma \ref{lem:compression} covers both mechanisms that occur in the manuscript with a single
statement, because both the renormalization maps $R^{N}_{N+M}$ and the passage from the full to the
chiral algebra are compressions along isometries: the latter is the inclusion of
$\fh^{(-)}_{N,\pm}=\ltwo(\Lambda_{N})_{\pm}\otimes p_{-}\C^{2}$ into
$\fh_{N,\pm}=\ltwo(\Lambda_{N})_{\pm}\otimes\C^{2}$ of (\mrChiral), with $p_{\pm}=\tfrac{1}{2}(\1_{2}\pm\sigma_{y})$.

\section{What the connecting unitaries actually require}
\label{sec:gns}

The manuscript's hypothesis is the existence of unitaries
$\phi^{N}_{M}:\cH^{(N)}_{M}\rightarrow\cH^{(N)}_{M+1}$ between the GNS spaces of $\omega^{(N)}_{M}$
and $\omega^{(N)}_{M+1}$, intertwining the two representations of $\fA_{N}$. Because $\fA_{N}$ is
finite-dimensional this is decided by a single integer.

\begin{prop}[GNS multiplicity]
\label{prop:gns}
Let $\fA\cong M_{d}(\C)$ and let $\omega$ be a state on $\fA$ with density matrix $\rho$. Then the
GNS representation of $\omega$ is unitarily equivalent to $\rank(\rho)$ copies of the unique
irreducible representation of $\fA$, so
\begin{align}
\label{eq:gnsdim}
\dim\cH_{\omega} & = d\cdot\rank(\rho),
\end{align}
and the GNS representations of two states $\omega,\omega'$ are unitarily equivalent if and only if
$\rank(\rho)=\rank(\rho')$. If moreover $\fA=\fA_{\CAR}(\fh)$ with $\dim\fh=d_{\fh}<\infty$ and
$\omega=\omega_{S}$ is quasi-free with symbol $S$, then
\begin{align}
\label{eq:qfrank}
\rank(\rho_{S}) & = 2^{m(S)}, & m(S) & := \#\{\text{eigenvalues of }S\text{ in }(0,1)\}
\end{align}
counted with multiplicity, so that $\dim\cH_{\omega_{S}}=2^{d_{\fh}}\cdot2^{m(S)}$.
\end{prop}

\begin{proof}
Every representation of $M_{d}(\C)$ is a multiple of the unique irreducible one on $\C^{d}$, and two
such are unitarily equivalent iff the multiplicities agree; the multiplicity of the GNS
representation of $\omega$ is $\rank(\rho)$, since $\cH_{\omega}=\fA\rho^{1/2}$ with the
Hilbert--Schmidt inner product, and $\fA\rho^{1/2}\cong\C^{d}\otimes\ran(\rho)$ as a left
$\fA$-module. For the quasi-free case, diagonalise $S$ with eigenvalues $\lambda_{1},\dots,
\lambda_{d_{\fh}}$; in the corresponding mode basis the density matrix factorises,
$\rho_{S}=\bigotimes_{i}\begin{psmallmatrix}\lambda_{i}&0\\0&1-\lambda_{i}\end{psmallmatrix}$, whose
rank is $\prod_{i}\rank\begin{psmallmatrix}\lambda_{i}&0\\0&1-\lambda_{i}\end{psmallmatrix}$; each
factor has rank $2$ if $\lambda_{i}\in(0,1)$ and rank $1$ if $\lambda_{i}\in\{0,1\}$. Finally
$\fA_{\CAR}(\fh)\cong M_{2^{d_{\fh}}}(\C)$.
\end{proof}

\begin{cor}[Existence of the Zini--Wang connecting unitaries]
\label{cor:phi}
In the situation of \eqref{eq:rgflow}--\eqref{eq:symbolflow}, an intertwining unitary
$\phi^{N}_{M}:\cH^{(N)}_{M}\rightarrow\cH^{(N)}_{M+1}$ exists if and only if
\begin{align}
\label{eq:phicrit}
m\big(S^{(N)}_{M}\big) & = m\big(S^{(N)}_{M+1}\big).
\end{align}
Purity of all the $\omega^{(N)}_{M}$, i.e.~$m\equiv0$, is sufficient for \eqref{eq:phicrit} but is
\emph{not} necessary.
\end{cor}

Corollary \ref{cor:phi} is the precise replacement for the manuscript's ``assuming a Stone--von
Neumann-type result at finite scales''. It also shows that the inference in the second quotation of
\S\ref{sec:intro} --- from non-purity of $\omega^{(N)}_{M}$ to the unavailability of the Zini--Wang
structure --- does not go through as stated: a family of mixed states of \emph{constant}
multiplicity supports connecting unitaries perfectly well. Section \ref{sec:obstruction} identifies
what does go wrong.

\section{The model of the manuscript}
\label{sec:model}

We now compute $m$ in each case. Recall from (\mrLatVac) the symbol of the massless lattice vacuum,
\begin{align}
\label{eq:latvac}
P(k) & = \tfrac{1}{2}\big(\1_{2}+\sign(k)\big(-\sin(\tfrac{1}{2}\vep_{N}k)\sigma_{x}
 +\cos(\tfrac{1}{2}\vep_{N}k)\sigma_{y}\big)\big), & k & \in\Gam_{N,\pm},
\end{align}
with $\sign(0)=0$, and from (\mrHardy) that the massless scaling-limit symbol is
$P_{0}(k)=\tfrac{1}{2}(\1_{2}+\sign(k)\sigma_{y})=p_{\sign(k)}$, i.e.~the Hardy projection. We write
$\vep_{N,M}:=\vep_{N+M}$ for the lattice spacing at the fine scale.

\begin{thm}[Closed forms]
\label{thm:model}
Let $\Gam_{N,-}$ be the Neveu--Schwarz momentum lattice, so that $0\notin\Gam_{N,-}$. Then:
\begin{enumerate}
\item[(i)] $P(k)$ of \eqref{eq:latvac} is a projection for every $k$, so the lattice vacuum is pure
on the \emph{full} algebra $\fA_{N,\pm}$;
\item[(ii)] with $p_{\mp}=\tfrac{1}{2}(\1_{2}\mp\sigma_{y})$,
\begin{align}
\label{eq:commutator}
\big\|[P(k),p_{\mp}]\big\| & = \tfrac{1}{2}\big|\sin(\tfrac{1}{2}\vep_{N}k)\big|,
\end{align}
which vanishes if and only if $\vep_{N}k\in2\pi\Z$;
\item[(iii)] the chiral compression $p_{-}P(k)p_{-}$ is the scalar
\begin{align}
\label{eq:chiraleig}
\lambda_{N}(k) & = \tfrac{1}{2}\big(1-\sign(k)\cos(\tfrac{1}{2}\vep_{N}k)\big), &
\lambda_{N}(1-\lambda_{N}) & = \tfrac{1}{4}\sin^{2}(\tfrac{1}{2}\vep_{N}k);
\end{align}
\item[(iv)] for the momentum-cutoff renormalization group of Definition \mrDefMomRG\ on the full algebra,
$S^{(N)}_{M}(k)=P(k)|_{\vep_{N}\rightarrow\vep_{N,M}}$ for $k\in\Gam_{N,\pm}$, which is a projection
for every $M\in\N_{0}\cup\{\infty\}$; hence $m(S^{(N)}_{M})=0$ throughout and
$\omega^{(N)}_{M}$ is pure;
\item[(v)] for the momentum-cutoff renormalization group on the \emph{chiral} algebra
$\fA^{(-)}_{N,-}$,
\begin{align}
\label{eq:mchiral}
m\big(S^{(N)}_{M}\big) & = |\Gam_{N,-}| \quad\text{for every finite }M, &
m\big(S^{(N)}_{\infty}\big) & = 0;
\end{align}
\item[(vi)] for the wavelet renormalization group of Definition \mrDefWavRG\ on the chiral algebra, applied
to the scaling-limit state, the symbol $(R^{N}_{\infty})^{*}P^{-}R^{N}_{\infty}$ is circulant on
$\ltwo(\Lambda_{N})$ with eigenvalues
\begin{align}
\label{eq:waveletspec}
\mu_{N}(\kappa) & = \!\!\!\sum_{j\in\Z\!\ :\!\ \kappa+2\pi j/\vep_{N}<0}\!\!\!
\big|\hat{s}(\vep_{N}\kappa+2\pi j)\big|^{2}, & \kappa & \in\Gam_{N,-},
\end{align}
and these satisfy $\mu_{N}(\kappa)\in(0,1)$ strictly, so $m=|\Gam_{N,-}|$, together with
\begin{align}
\label{eq:waveletbound}
\min\{\mu_{N}(\kappa),1-\mu_{N}(\kappa)\} & \leq \delta(\vep_{N}\kappa)
 := 1-\big|\hat{s}(\vep_{N}\kappa)\big|^{2} = \mathcal{O}\big((\vep_{N}\kappa)^{2K}\big);
\end{align}
\item[(vii)] on the chiral algebra, for \emph{either} renormalization group and every finite $M$, one has $m(S^{(N)}_{M})=|\Gam_{N,-}|$.
\end{enumerate}
\end{thm}

\begin{proof}
(i) The vector $\vec{n}(k)=\sign(k)(-\sin(\tfrac{1}{2}\vep_{N}k),\cos(\tfrac{1}{2}\vep_{N}k),0)$ is a
unit vector for $k\neq0$, and $\tfrac{1}{2}(\1_{2}+\vec{n}\cdot\vec{\sigma})$ is a rank-one
projection; for $k=0$ (Ramond only) $P(0)=\tfrac{1}{2}\1_{2}$, cf.~the text following (\mrLatVac).

(ii) Only the $\sigma_{x}$-part fails to commute with $\sigma_{y}$:
\begin{align*}
[P(k),p_{\mp}] & = \mp\tfrac{1}{2}\cdot\tfrac{1}{2}\sign(k)\big(-\sin(\tfrac{1}{2}\vep_{N}k)\big)
 [\sigma_{x},\sigma_{y}] = \pm\tfrac{i}{2}\sign(k)\sin(\tfrac{1}{2}\vep_{N}k)\!\ \sigma_{z},
\end{align*}
using $[\sigma_{x},\sigma_{y}]=2i\sigma_{z}$ and $[\sigma_{y},\sigma_{y}]=0$. Since $\|\sigma_{z}\|=1$
and $|\sign(k)|=1$ for $k\neq0$, \eqref{eq:commutator} follows.

(iii) With $e_{\pm}=2^{-1/2}(e_{1}\pm ie_{2})$ the eigenvectors of $\sigma_{y}$, one has
$\sigma_{y}e_{-}=-e_{-}$ and $\sigma_{x}e_{-}=-ie_{+}\perp e_{-}$, so
$\langle e_{-},P(k)e_{-}\rangle=\tfrac{1}{2}(1-\sign(k)\cos(\tfrac{1}{2}\vep_{N}k))$. The second
identity in \eqref{eq:chiraleig} is elementary. Consistency with Lemma \ref{lem:compression} is the
computation $\|X\|^{2}=\|[P,p_{-}]\|^{2}=\tfrac{1}{4}\sin^{2}(\tfrac{1}{2}\vep_{N}k)
=\lambda_{N}(1-\lambda_{N})=\|S'-S'^{2}\|$.

(iv) By (\mrMomRG) the momentum-cutoff isometry acts as multiplication by
$\vep^{1/2}\chi_{\Gam_{N}}$ on the Fourier side, so $\ran R^{N}_{N+M}$ is the span of the momentum
modes in $\Gam_{N}\subset\Gam_{N+M}$, with both $\C^{2}$-components. The symbol $S^{(N+M)}$ acts
fibrewise in $k$, hence leaves this span invariant; Lemma \ref{lem:compression}(iii) applies and the
compression is the restriction of $P$ to those fibres, still a projection. Letting $M\rightarrow
\infty$ replaces $P$ by $P_{0}$, again a projection.

(v) By (iii) the compressed chiral symbol at $k\in\Gam_{N,-}$ is the scalar $\lambda_{N+M}(k)$ of
\eqref{eq:chiraleig} with $\vep_{N}$ replaced by $\vep_{N,M}$. It lies in $(0,1)$ unless
$\sin(\tfrac{1}{2}\vep_{N,M}k)=0$, i.e.~unless $\vep_{N,M}k\in2\pi\Z$. For $k\in\Gam_{N,-}$ we have
$|\vep_{N,M}k|\leq\vep_{N,M}\pi/\vep_{N}=2^{-M}\pi<2\pi$, and $\vep_{N,M}k=0$ is excluded because
$0\notin\Gam_{N,-}$. Hence every one of the $|\Gam_{N,-}|$ eigenvalues lies strictly in $(0,1)$,
for every finite $M$. As $M\rightarrow\infty$, $\lambda_{N+M}(k)\rightarrow
\tfrac{1}{2}(1-\sign(k))=\chi_{k<0}\in\{0,1\}$, so $m$ drops to $0$.

(vi) By Definition \mrDefWavRG\ the wavelet isometry acts on the Fourier side as multiplication by
$\vep_{N}^{1/2}\hat{s}(\vep_{N}\!\ \cdot\!\ )$ followed by periodisation, so for
$\xi,\eta\in\ltwo(\Lambda_{N})$
\begin{align*}
\big\langle\xi,(R^{N}_{\infty})^{*}P^{-}R^{N}_{\infty}\eta\big\rangle
 & = \frac{1}{|\Lambda_{N}|}\sum_{l\in\Gam_{\infty,-},\!\ l<0}\big|\hat{s}(\vep_{N}l)\big|^{2}
 \overline{\hat{\xi}_{l}}\!\ \hat{\eta}_{l},
\end{align*}
which is invariant under simultaneous translation of $\xi,\eta$ on $\Lambda_{N}$, hence circulant;
its eigenvalue at the lattice momentum $\kappa\in\Gam_{N,-}$ is the sum of $|\hat{s}|^{2}$ over the
negative members of the fibre $\kappa+\tfrac{2\pi}{\vep_{N}}\Z$, which is \eqref{eq:waveletspec}.
Orthonormality of $\{s(\cdot-n)\}_{n\in\Z}$ is equivalent to $\sum_{j\in\Z}|\hat{s}(\xi+2\pi j)|^{2}
=1$, so the full fibre sums to $1$ and
\begin{align*}
1-\mu_{N}(\kappa) & = \!\!\!\sum_{j\!\ :\!\ \kappa+2\pi j/\vep_{N}>0}\!\!\!
 \big|\hat{s}(\vep_{N}\kappa+2\pi j)\big|^{2}.
\end{align*}
If $\kappa<0$ the term $j=0$ occurs in $\mu_{N}$, whence $1-\mu_{N}(\kappa)\leq\sum_{j\neq0}
|\hat{s}(\vep_{N}\kappa+2\pi j)|^{2}=1-|\hat{s}(\vep_{N}\kappa)|^{2}=\delta(\vep_{N}\kappa)$; if
$\kappa>0$ the same argument bounds $\mu_{N}(\kappa)$. This is \eqref{eq:waveletbound}, and
$\delta(\xi)=\mathcal{O}(\xi^{2K})$ is (\mrAlias) of the manuscript. Strict positivity of both
$\mu_{N}$ and $1-\mu_{N}$ holds because $\hat{s}$ vanishes only on $2\pi\Z\setminus\{0\}$ while
$\vep_{N}\kappa+2\pi j\in2\pi\Z$ would force $\vep_{N}\kappa\in2\pi\Z$, excluded as in (v). Indeed
the single term $j=-1$ (for $\kappa>0$) gives the matching lower bound
$\mu_{N}(\kappa)\geq|\hat{s}(\vep_{N}\kappa-2\pi)|^{2}\geq c\!\ (\vep_{N}\kappa)^{2K}$, since
$\hat{s}$ has a zero of order exactly $K$ at $-2\pi$.

(vii) Both renormalization groups act on the Fourier side by multiplication with a \emph{scalar} symbol --- $\vep^{1/2}\chi_{\Gam_{N}}$ in the momentum-cutoff case (\mrMomRG), $\vep_{N}^{1/2}\hat{s}(\vep_{N}\!\ \cdot\!\ )$ in the wavelet case --- hence commute with the constant spinor projections $p_{\pm}$ and restrict to the chiral one-particle space. By (iii) and (v) the chiral symbol at the fine scale is the multiplication operator $\lambda_{N+M}$, which satisfies $0<\lambda_{N+M}<1$ strictly on $\Gam_{N+M,-}$; \eqref{eq:strictm} of Lemma \ref{lem:compression} therefore gives $m(S^{(N)}_{M})=\dim\fh^{(-)}_{N,-}=|\Gam_{N,-}|$, whichever isometry is used.
\end{proof}

\begin{rem}[The aliasing mass is the same object twice]
\label{rem:aliasing}
The quantity $\delta(\xi)=1-|\hat{s}(\xi)|^{2}=\sum_{j\neq0}|\hat{s}(\xi-2\pi j)|^{2}$ controlling
\eqref{eq:waveletbound} is exactly the aliasing mass that drives the Jackson estimate of
Lemma \mrJackson\ in the manuscript, and its order $2K$ is (\mrAlias). The failure of the wavelet
renormalization group to produce pure finitely renormalized states is therefore the \emph{same}
phenomenon as the failure of $P^{\pm}$ to preserve the wavelet core $\cD_{W}$ recorded in
Remark \mrMoeb: both say that a multiresolution subspace is not a spectral subspace of a Fourier
multiplier, and both are quantified by $\delta$.
\end{rem}

\section{The obstruction, and its size}
\label{sec:obstruction}

\begin{cor}[Where the Zini--Wang structure really fails]
\label{cor:obstruction}
For the chiral algebra $\fA^{(-)}_{N,-}$ with the momentum-cutoff renormalization group,
\begin{align}
\label{eq:dimcollapse}
\dim\cH^{(N)}_{M} & = 2^{|\Gam_{N,-}|}\cdot2^{|\Gam_{N,-}|}\ \ \text{for every finite }M,
 & \dim\cH^{(N)}_{\infty} & = 2^{|\Gam_{N,-}|}.
\end{align}
Consequently:
\begin{enumerate}
\item[(i)] the connecting unitaries $\phi^{N}_{M}:\cH^{(N)}_{M}\rightarrow\cH^{(N)}_{M+1}$ of the
manuscript \emph{do} exist, for every finite $M$, by Corollary \ref{cor:phi} --- notwithstanding
that every $\omega^{(N)}_{M}$ is mixed;
\item[(ii)] but no such unitary exists between any $\cH^{(N)}_{M}$ and $\cH^{(N)}_{\infty}$: the GNS
multiplicity collapses by the factor $2^{|\Gam_{N,-}|}$ in the limit.
\end{enumerate}
Hence the inductive limit of the $\cH^{(N)}_{M}$ along the $\phi^{N}_{M}$ is \emph{not} the
scaling-limit representation. Two comparisons are worth drawing. On the full algebra $\fA_{N,\pm}$, Theorem \ref{thm:model}(iv) gives $m\equiv0$ throughout and no collapse occurs. On the chiral algebra with the \emph{wavelet} renormalization group, Theorem \ref{thm:model}(vi),(vii) gives $m=|\Gam_{N,-}|$ at every $M$ \emph{including} $M=\infty$, so again no collapse occurs: the wavelet route keeps a mixed limit state and with it a constant multiplicity. The collapse is thus specific to the momentum-cutoff route, whose limit state is pure precisely because its isometries are spectral projections of the same multiplier that defines the Hardy projection.
\end{cor}

\begin{proof}
Immediate from Proposition \ref{prop:gns} and Theorem \ref{thm:model}(iv),(v), with
$d_{\fh}=|\Gam_{N,-}|$ for the chiral one-particle space.
\end{proof}

This is the precise form of the manuscript's assertion, and it differs from it in an instructive
way. The obstruction is not that the finitely renormalized states fail to be pure --- they do fail,
but that alone is harmless. The obstruction is that they fail to be pure \emph{by an amount that does
not survive the limit}: the multiplicity is locally constant in $M$ and then jumps at $M=\infty$.
The manuscript's parenthetical ``(due to the entanglement of the chiral halves)'' is vindicated as
the correct \emph{mechanism} --- by Theorem \ref{thm:model}(ii) the chiral halves are entangled
exactly to the extent that $[P(k),p_{\mp}]\neq0$ --- while ``not pure'' is the wrong \emph{criterion}.

\begin{cor}[The mixing is extensive, with density $2\ln2-1$]
\label{cor:entropy}
The chiral restriction of the finite-scale massless lattice vacuum has von Neumann entropy
\begin{align}
\label{eq:entropy}
S\big(\omega^{(N)}_{0}|_{\fA^{(-)}_{N,-}}\big) & = \sum_{k\in\Gam_{N,-}}H_{2}\big(\lambda_{N}(k)\big),
 & H_{2}(p) & := -p\ln p-(1-p)\ln(1-p),
\end{align}
with $\lambda_{N}$ as in \eqref{eq:chiraleig}, and the entropy density converges:
\begin{align}
\label{eq:entropydensity}
\lim_{N\rightarrow\infty}\frac{1}{|\Gam_{N,-}|}
 S\big(\omega^{(N)}_{0}|_{\fA^{(-)}_{N,-}}\big)
 & = \frac{1}{2\pi}\int_{-\pi}^{\pi}\!H_{2}\Big(\frac{1-\cos(u/2)}{2}\Big)du = 2\ln2-1
 \approx0.386294.
\end{align}
The mixing is concentrated at the Brillouin-zone boundary, where
$\lambda_{N}(1-\lambda_{N})\rightarrow\tfrac{1}{4}$ is maximal, and vanishes as
$\tfrac{1}{4}(\tfrac{1}{2}\vep_{N}k)^{2}+\mathcal{O}((\vep_{N}k)^{4})$ on the modes that survive the
scaling limit.
\end{cor}

\begin{proof}
For a quasi-free state the entropy is $\sum_{i}H_{2}(\lambda_{i})$ over the eigenvalues of the
symbol, which with \eqref{eq:chiraleig} is \eqref{eq:entropy}. Writing $u=\vep_{N}k$, the momenta
$k\in\Gam_{N,-}$ equidistribute in $u\in(-\pi,\pi)$ as $N\rightarrow\infty$, and
$\lambda_{N}=\tfrac{1}{2}(1-\sign(u)\cos(u/2))$, whose binary entropy is even in $u$ and equals
$H_{2}(\tfrac{1}{2}(1-\cos(u/2)))$; this gives the Riemann sum in \eqref{eq:entropydensity}. For the
value, substitute $u=2t$ and then $t=2\varphi$, and use
$H_{2}(\sin^{2}\varphi)=-2\sin^{2}\varphi\ln\sin\varphi-2\cos^{2}\varphi\ln\cos\varphi$:
\begin{align*}
\frac{1}{2\pi}\int_{-\pi}^{\pi}\!H_{2}\Big(\frac{1-\cos(u/2)}{2}\Big)du
 & = \frac{2}{\pi}\int_{0}^{\pi/2}\!H_{2}\big(\sin^{2}(t/2)\big)dt
 = -\frac{8}{\pi}\int_{0}^{\pi/4}\!\big[\sin^{2}\!\varphi\ln\sin\varphi
 +\cos^{2}\!\varphi\ln\cos\varphi\big]d\varphi.
\end{align*}
The substitution $\varphi\mapsto\tfrac{\pi}{2}-\varphi$ turns the second summand into the first over
$[\tfrac{\pi}{4},\tfrac{\pi}{2}]$, so the bracket integrates to
$\int_{0}^{\pi/2}\sin^{2}\!\varphi\ln\sin\varphi\!\ d\varphi=\tfrac{\pi}{8}(1-\ln4)$, whence the
value $-\tfrac{8}{\pi}\cdot\tfrac{\pi}{8}(1-2\ln2)=2\ln2-1$.
\end{proof}

So the chiral restriction is not merely mixed: it is mixed at maximal rate on a finite fraction of
the modes, and its entropy grows like $(2\ln2-1)|\Gam_{N,-}|$. What makes the scaling limit pure
nevertheless is that the strongly mixed modes are precisely the lattice artefacts near
$|k|\sim\pi/\vep_{N}$, which are pushed to infinity by the limit, while each fixed mode is mixed only
to order $(\vep_{N}k)^{2}$.

\section{What survives unconditionally}
\label{sec:survives}

For completeness we record the part of the comparison that requires no hypotheses, and which the
manuscript states correctly.

\begin{prop}[OAR supplies Zini--Wang-compatible data]
\label{prop:oarzw}
Let $\{\fA_{N},\alpha^{N}_{N+1}\}$ be an OAR inductive system with a projectively consistent family
of scaling-limit states $\omega^{(N)}_{\infty}$, and let
$\{\cH^{(N)}_{\infty},\pi^{(N)}_{\infty},\Omega^{(N)}_{\infty}\}$ be the GNS data. Then there are
isometries $V^{N}_{N+1}:\cH^{(N)}_{\infty}\rightarrow\cH^{(N+1)}_{\infty}$ with
$V^{N}_{N+1}\Omega^{(N)}_{\infty}=\Omega^{(N+1)}_{\infty}$ and
\begin{align}
\label{eq:intertwine}
V^{N}_{N+1}\!\ \pi^{(N)}_{\infty}(O) & = \pi^{(N+1)}_{\infty}\big(\alpha^{N}_{N+1}(O)\big)V^{N}_{N+1},
 & O & \in\fA_{N},
\end{align}
and consequently
\begin{align}
\label{eq:proj}
V^{N}_{N+1}\!\ \pi^{(N)}_{\infty}(O)\!\ (V^{N}_{N+1})^{*}
 & = \pi^{(N+1)}_{\infty}\big(\alpha^{N}_{N+1}(O)\big)\!\ p_{\ran V^{N}_{N+1}}.
\end{align}
Moreover $V^{N}_{N+1}$ is unitary if and only if
$\pi^{(N+1)}_{\infty}(\alpha^{N}_{N+1}(\fA_{N}))\Omega^{(N+1)}_{\infty}$ is dense in
$\cH^{(N+1)}_{\infty}$. In the quasi-free case with $\omega^{(N)}_{\infty}$ pure, $V^{N}_{N+1}$ is
the second quantisation $\Gamma(R^{N}_{N+1})$ of the one-particle isometry, which is exactly the
shape of a Zini--Wang \emph{strong} scaling limit.
\end{prop}

\begin{proof}
Projective consistency $\omega^{(N+1)}_{\infty}\circ\alpha^{N}_{N+1}=\omega^{(N)}_{\infty}$ means
that $\pi^{(N)}_{\infty}(O)\Omega^{(N)}_{\infty}\mapsto
\pi^{(N+1)}_{\infty}(\alpha^{N}_{N+1}(O))\Omega^{(N+1)}_{\infty}$ preserves inner products, hence
extends to an isometry $V^{N}_{N+1}$ on the closure, which is \eqref{eq:intertwine} by construction;
\eqref{eq:proj} follows by multiplying \eqref{eq:intertwine} with $(V^{N}_{N+1})^{*}$ on the right
and using $V^{*}V=\1$, $VV^{*}=p_{\ran V}$. The range of $V^{N}_{N+1}$ is by construction the closure
of $\pi^{(N+1)}_{\infty}(\alpha^{N}_{N+1}(\fA_{N}))\Omega^{(N+1)}_{\infty}$, giving the unitarity
criterion. The last statement is the standard identification of the GNS representation of a pure
quasi-free state with the Fock representation over the one-particle space.
\end{proof}

Equation \eqref{eq:proj} is the identity displayed in \S1.2 of the manuscript; Proposition
\ref{prop:oarzw} supplies its proof and the accompanying unitarity criterion.

\section{Summary}
\label{sec:summary}

\begin{center}
\small
\begin{tabular}{lcccc}
\toprule
 & $m(S^{(N)}_{M})$, $M<\infty$ & $m(S^{(N)}_{\infty})$ & $\phi^{N}_{M}$ exist? & reaches the limit? \\
\midrule
full algebra, mom.\ cutoff & $0$ & $0$ & yes & yes \\
chiral algebra, mom.\ cutoff & $|\Gam_{N,-}|$ & $0$ & yes & \textbf{no} \\
chiral algebra, wavelet & $|\Gam_{N,-}|$ & $|\Gam_{N,-}|$ & yes & yes \\
\bottomrule
\end{tabular}
\end{center}

\noindent Only the first row consists of pure states; rows two and three are mixed throughout, and mixedness is visibly not what decides the last column.

\medskip

\noindent The three conclusions, in order of importance:
\begin{enumerate}
\item[(1)] The criterion for the Zini--Wang connecting unitaries is constancy of the GNS
multiplicity $m$, not purity (Corollary \ref{cor:phi}). The manuscript's hypothesis ``assuming a
Stone--von Neumann-type result at finite scales'' is exactly \eqref{eq:phicrit}.
\item[(2)] For the chiral algebra the multiplicity is constant at every finite scale and collapses
only in the limit, so the obstruction is a discontinuity at $M=\infty$, not the mixedness of the
$\omega^{(N)}_{M}$ (Corollary \ref{cor:obstruction}); and the collapse is specific to the momentum-cutoff route, the wavelet route having constant multiplicity all the way to $M=\infty$. The mechanism named in the manuscript --- the
entanglement of the chiral halves --- is correct, and is quantified by
$\|[P(k),p_{\mp}]\|=\tfrac{1}{2}|\sin(\tfrac{1}{2}\vep_{N}k)|$.
\item[(3)] The failure of purity is extensive, with entropy density exactly $2\ln2-1$ per site
(Corollary \ref{cor:entropy}); it is carried by the zone-boundary lattice artefacts and is of order
$(\vep_{N}k)^{2}$ on the modes that survive.
\end{enumerate}

\section{What this note does not settle}
\label{sec:open}

\begin{itemize}
\item \emph{The converse direction.} We have shown that OAR supplies Zini--Wang data (Proposition
\ref{prop:oarzw}) and determined exactly when the connecting unitaries exist. We have not shown that
a Zini--Wang low-energy scaling limit conversely determines an OAR inductive system; this would
require reconstructing the refining morphisms $\alpha$ from the $\phi$'s, and the identity
\eqref{eq:proj} shows that the $\phi$'s see only the compression of $\pi(\fA_{N})$ to $\ran V$.
\item \emph{Beyond quasi-free.} Proposition \ref{prop:gns} holds for arbitrary states on a
finite-dimensional algebra, but the closed forms of Theorem \ref{thm:model} use quasi-freeness
throughout. The Zini--Wang conjecture that higher-level anyonic chains admit compatible embeddings
\cite[p.~898]{ZiniWang} is untouched by these methods.
\item \emph{Infinite scales.} All statements are at fixed finite $N$, where $\fA_{N}$ is a matrix
algebra. The passage $N\rightarrow\infty$ leaves the setting of Proposition \ref{prop:gns}, and the
relevant criterion there is quasi-equivalence of quasi-free states, governed by the Powers--St\o rmer
condition that $S^{1/2}-S'^{1/2}$ be Hilbert--Schmidt \cite{Araki,PowersStormer}. We have not
carried out that analysis.
\item \emph{The strong scaling limit.} Zini and Wang's strong scaling limit requires the $\phi$'s to
come from a single isometry $\phi_{N}:\cH_{N}\rightarrow\cH_{N+1}$ of the full spaces. Proposition
\ref{prop:oarzw} produces such an isometry in the pure quasi-free case, but we have not examined
whether the mixed case of Theorem \ref{thm:model}(v)--(vii) admits one.
\end{itemize}

\section*{Numerical verification}

All closed forms above are checked in \texttt{numerics/os\_check19.py} (pure \texttt{numpy}):
\begin{itemize}
\item \texttt{[31]} $\|P^{2}-P\|\leq1.1\cdot10^{-16}$; \eqref{eq:commutator} to machine zero
($0.0$); \eqref{eq:chiraleig} to $3.3\cdot10^{-16}$; the entropy density of \eqref{eq:entropydensity}
agrees with $2\ln2-1$ to $10$ decimal places.
\item \texttt{[32]} the momentum-cutoff compression satisfies $\|S'^{2}-S'\|\leq7.9\cdot10^{-17}$ for
all tested $N,M$ --- Theorem \ref{thm:model}(iv).
\item \texttt{[33]} the circulant spectrum \eqref{eq:waveletspec} reproduces the directly computed
eigenvalues to $10^{-15}$ (db4, db6) and the bound \eqref{eq:waveletbound} holds at every $\kappa$.
For db2 the agreement is only $10^{-10}$, a replica-truncation artefact. Note that
$\min\{\mu,1-\mu\}$ can be as small as $10^{-12}$, so the strict positivity asserted in Theorem
\ref{thm:model}(vi) is at the edge of double precision and is established analytically, not
numerically.
\item \texttt{[34]} the multiplicities of \eqref{eq:mchiral}: $m=0$ throughout on the full algebra;
$m=|\Gam_{N,-}|$ for $M=1,\dots,7$ on the chiral algebra and $m=0$ at $M=\infty$, for
$N=3,5$.
\end{itemize}
A numerically stable evaluation of $\hat{s}$ is required for \texttt{[33]}: the order-$K$ zero of
$m_{0}$ at $\pi$ is deflated in closed form, since evaluating $m_{0}(\pi+\eta)\sim\eta^{K}$ from the
filter coefficients loses all precision.

\begin{thebibliography}{9}

\bibitem{Araki}
H.~Araki.
\newblock On quasifree states of CAR and Bogoliubov automorphisms.
\newblock \emph{Publ. Res. Inst. Math. Sci.} \textbf{6} (1970/71), 385--442.

\bibitem{EvansKawahigashi}
D.~E. Evans and Y.~Kawahigashi.
\newblock \emph{Quantum Symmetries on Operator Algebras}.
\newblock Oxford University Press, 1998.

\bibitem{OsborneStottmeister}
T.~J. Osborne and A.~Stottmeister.
\newblock Conformal field theory from lattice fermions.
\newblock \emph{Comm. Math. Phys.} \textbf{398} (2023), 219--289; arXiv:2107.13834.
\newblock Referenced here in the revised form \texttt{free\_fermion\_cft\_v5.tex}.

\bibitem{PowersStormer}
R.~T. Powers and E.~St\o rmer.
\newblock Free states of the canonical anticommutation relations.
\newblock \emph{Comm. Math. Phys.} \textbf{16} (1970), 1--33.

\bibitem{ZiniWang}
M.~S. Zini and Z.~Wang.
\newblock Conformal field theories as scaling limit of anyonic chains.
\newblock \emph{Comm. Math. Phys.} \textbf{363} (2018), 877--953.

\end{thebibliography}

\end{document}
